\section{The broad pure-quadratic reflection reduction on the same window}
\label{sec:abelian-reflection-actual-window}

\paragraph{Status and exact scope.}
This conditional-interface proof applies the original-coefficient
strict theorem AC-G168 to the four actual finite labels
$1,\chi_{-4},\chi_8,\chi_{-8}$ over $K$.
It retains every omission of AC-G167-window.
The old original-coefficient residue-zero theorem is used only where
it was already proved; residual zero in AC-G168 is not promoted
silently to an integral assertion.
One full local comparison cost covers the entire enlarged union.
The remaining numerical bound and mixed realization stay open.

Put $p=269$, $K=\mathbb Q(\omega)$, and $t=\sqrt[3]p>0$.
Let $\chi:G_K\to\mu_3$ describe $K(t)/K$, let $\phi$ be the
primitive odd quartic character modulo sixteen with $\phi(13)=i$
and kernel $\{1,7\}$, and put $\delta_d=\chi_d|_{G_K}$ for
$d=1,-4,8,-8$, with $\delta_1=1$.
For the retained datum, $0<\varepsilon<1/256$, page and
prescribed pair $q,d$, keep the actual objects
\[
 \begin{aligned}
 F_u&=\mathbb Q(\zeta_{2^{h(u)+2}},\zeta_7,
                                  \sqrt{-q},\sqrt{-d},\omega),\\
 H_u&=F_u(t),\quad G_u=\operatorname{Gal}(H_u/K),\quad D_u=|G_u|,\\
 u^{2+8\varepsilon}&\leq2^{h(u)}\leq2u^{2+8\varepsilon},
                      \qquad T=\lfloor u^{399}\rfloor .
 \end{aligned}
 \tag{ARW.1}
\]
The prescribed pair's $d$ is separate from the subscript of
$\delta_d$.  All original $O_\ell,E_\ell,b_\ell$, strict local
conditions and derived indices retain (FC.1)--(FC.2), (EU.2)
and (TC.1).  Every coefficient prime remains in that source range.

Write
\[
 \mathcal E=\{89,92,99,168,175,178\}\subset\mathbb Z/268\mathbb Z.
 \tag{ARW.2}
\]
For $u\geq2$ define the unchanged part of the omitted union by
\[
 \begin{aligned}
 Z_{16}^+&=\{(p,\chi^{-j}\phi^{-\sigma},e):j=1,2,\ \sigma=\pm1,\\
 &\hspace{2em}1\leq e<T,\ e\equiv-1\pmod{67}\},\\
 Z_{\rm quad}^+&=\{(p,\chi^{-j}\delta_d,e):j=1,2,\ d=-4,8,-8,\\
 &\hspace{2em}1\leq e<T,\ e\equiv-1\pmod{67}\},\\
 Z_{\rm cub}^{\ne0}&=\{(p,\chi^j,e):j=1,2,\ 1\leq e<T,\\
 &\hspace{2em}e\equiv0\pmod{268}\ \text{or}\ e\equiv-1\pmod{67}\},\\
 Z_{16}^{0,q}&=\{(p,\phi^{-\sigma},e):\sigma=\pm1,\ 1\leq e<T,\\
 &\hspace{2em}e\equiv66,200\pmod{268}\},\\
 Z_3&=\{(3,1,e):1\leq e<T\},\\
 Z_{\rm keep}&=Z_{16}^+\sqcup Z_{\rm quad}^+
              \sqcup Z_{\rm cub}^{\ne0}\sqcup Z_{16}^{0,q}\sqcup Z_3.
 \end{aligned}
 \tag{ARW.3}
\]
The enlarged pure-label part is
\[
 \begin{aligned}
 Z_1^{\rm all}&=\{(p,1,e):1\leq e<T\},\\
 Z_d^{\rm ref}&=\{(p,\delta_d,e):1\leq e<T,\ e\not\equiv0\pmod{268}\}
                                  \quad(d=-4,8),\\
 Z_{-8}^{\rm ref}&=\{(p,\delta_{-8},e):1\leq e<T,
                        \ e\bmod268\notin\{0\}\cup\mathcal E\},\\
 Z^{\rm ref}&=Z_{\rm keep}\sqcup Z_1^{\rm all}
              \sqcup Z_{-4}^{\rm ref}\sqcup Z_8^{\rm ref}\sqcup Z_{-8}^{\rm ref}.
 \end{aligned}
 \tag{ARW.4}
\]
These unions are disjoint.  Let $I_{399}^{\rm estreflection}$ be
(TC.1) with exactly $Z^{\rm ref}$ omitted.  The predecessor is
$I_{399}^{\rm estpure16quarter}$ of AC-G167-window, whose omitted
set $Z^{0,16q}$ is (P16QW.2).

For positive integer $T$ use the exact counting functions
\[
 \begin{aligned}
 N_T&=T-1,& B_T&=\lfloor(T-1)/268\rfloor,\\
 A_T&=\lfloor T/67\rfloor,& Q_T&=\lfloor(T+67)/134\rfloor,\\
 R_T&=\sum_{r\in\mathcal E}\left\lfloor\frac{T+267-r}{268}\right\rfloor.
 \end{aligned}
 \tag{ARW.5}
\]

\begin{theorem}[AC-G169-window: the broad reflection reduction]
\label{AC-G169-window}
Every triple in (ARW.3)--(ARW.4) is an actual host character,
coefficient prime and weight whose original strict group is zero.
Every old AC-G167 omission is retained.  The exact counts are
\[
 \begin{aligned}
 |Z_{16}^+|&=4A_T,& |Z_{\rm quad}^+|&=6A_T,\\
 |Z_{\rm cub}^{\ne0}|&=2(B_T+A_T),&|Z_{16}^{0,q}|&=2Q_T,\\
 |Z_3|&=N_T,&|Z_1^{\rm all}|&=N_T,\\
 |Z_{-4}^{\rm ref}|&=N_T-B_T,&|Z_8^{\rm ref}|&=N_T-B_T,\\
 |Z_{-8}^{\rm ref}|&=N_T-B_T-R_T.
 \end{aligned}
 \tag{ARW.6}
\]
In particular the total coefficient-$269$ count and the number
of newly omitted triples relative to AC-G167 are respectively
\[
 \begin{aligned}
 |Z^{\rm ref}_{269}|&=4N_T-B_T+12A_T+2Q_T-R_T,\\
 |Z^{\rm ref}\setminus Z^{0,16q}|&=4(N_T-B_T-A_T)-R_T.
 \end{aligned}
 \tag{ARW.7}
\]
For $T\geq2$ the latter set is nonempty.
No original-coefficient residue-zero assertion for the three
nontrivial quadratic labels, or zero at the six uncovered
$\delta_{-8}$ residues, is inferred.

Let $C_u$ be the actual normalized positive-index cost over
this entire omitted union.  In nonnegative extended sums,
\[
 \begin{aligned}
 I_{399}&=I_{399}^{\rm estreflection}+C_u,\\
 I_{399}^{\rm estreflection}&\leq I_{399}^{\rm estpure16quarter}\leq I_{399},\\
 0\leq C_u&\leq8D_u^2T^2\log(8TD_u)
                +\tfrac12D_u\mathcal R_uT^2+2D_ur_uc_0T\\
 &=O_\varepsilon\bigl(u^{802+16\varepsilon}(\log u)^2\bigr).
 \end{aligned}
 \tag{ARW.8}
\]
One full (FL.4) majorant covers the entire old and new union;
no individual positive index is asserted zero.
The original common-window premise (TC.3), or its AC-G167
version, may use $I_{399}^{\rm estreflection}$ without changing
any simultaneous field, coefficient, finite-character or weight
quantifier.  An older premise implies the new one by nonnegativity.
Conversely a new positive-relative-slack premise implies the
full original premise on the same $W,J$, datum, page and pair,
after increasing its initial threshold and halving slack.
The actual $R.t$ and every retained clause of
$\mathcal R_{\rm cub}$ remain required and unproved.
\end{theorem}

\begin{proof}
\proofstep{1. Actual labels, full local conditions and nonzero residues.}
The prescribed host contains $K(t)$ and $\mathbb Q(\zeta_{16})$.
Its abelian factor is disjoint from $K(t)$ over $K$.
Moreover $K\cap\mathbb Q(\zeta_{16})=\mathbb Q$.
Thus the four pure quadratic labels $1,\delta_{-4},\delta_8,
\delta_{-8}$ are distinct, and differ from the two pure quartic
labels and every label with nontrivial cubic part.
The coefficient-three row is disjoint from all coefficient-$269$
rows.  These facts prove the disjointness of (ARW.3)--(ARW.4).

For the four pure quadratic labels the actual coefficient is
\[
 \begin{aligned}
 A_{d,e}&=(O_p\otimes_{\mathbb Z_p}\mathbb Q_p/\mathbb Z_p)
                                                  (\delta_d)(-e)\\
        &=(E_p/O_p)(\delta_dc_p^{-e}).
 \end{aligned}
 \tag{ARW.9}
\]
For every nonzero residue retained in (ARW.4), apply (G168.3)
with this $d$, $m=0$ and $\kappa=1$.
This is exactly the original coefficient (ARW.9), with no
restriction on the absolute ramification of $O_p$.
Its strict classes restrict to zero in full $H^1(G_{K_v},A_{d,e})$
at $v\mid p$ and in $H^1(I_v,A_{d,e})$ at every other finite place.
The auxiliary field $\mathbb Q(\zeta_{24p})$ in its proof
is not a new original host or a new finite character label.

These conditions agree with (FC.2), including its ramification set.
Indeed (ARW.9) is unramified outside $2,3,p$.
At an added away-$p$ prime with trivial inertia action, a
coboundary restriction is a zero cocycle.  The cocycle identity
makes it zero on conjugate inertia groups and their closed
normal closure.  It descends to the smaller ramification set;
conversely inflation is injective and preserves the other
full local conditions.  Thus the global strict theorem and
the actual strict kernel in (FC.2) agree for the present use.
Every strict zero in $Z_{\rm keep}$ is retained from AC-G167,
including the coefficient-three statement without averaging by three.

\proofstep{2. Exactly the old original-coefficient residue-zero input.}
For $e\equiv0\pmod{268}$ the original strict zero of the
trivial label is already part of AC-G167.
For completeness its independent source is (K31.2), transported
in the complete proof of (EW.5)--(EW.6) of AC-G146.
Put $E_{\rm c}=K(t)$.  That input states
$S_p(E_{\rm c},e)=0$ with coefficient
$\mathbb Q_p/\mathbb Z_p(-e)$, unramified outside $p$ and
zero in full local cohomology at the unique $p$-place.
Its local invariant correction was included in the K31 proof.

Restrict an original strict class for $(E_p/O_p)c_p^{-e}$
from $K$ to $E_{\rm c}$.
Full local zero and inertia zero restrict to zero upstairs.
After removing added unramified primes, its target is
$O_p\otimes_{\mathbb Z_p}S_p(E_{\rm c},e)=0$.
This tensor identity uses a finite free $\mathbb Z_p$-basis
of the same $O_p$ in global and local continuous cochains.
Corestriction after restriction is multiplication by three,
which is invertible on the original coefficient; hence the
original class was zero.  Together with Step~1 this proves
$Z_1^{\rm all}$ for every positive weight.

For $j=1,2$ the same already-proved residue-zero rows remain
inside $Z_{\rm cub}^{\ne0}$; no new weight for those labels
is inferred from AC-G168.
For $\delta_{-4},\delta_8,\delta_{-8}$, the theorem's residual
zero at $r=0$ is not used to obtain an original strict zero.
Their local fixed-module obstruction at that residue remains
as stated in AC-G168.  Multiples of $268$ are therefore retained
in the unknown sum for all three nontrivial pure quadratic labels.

\proofstep{3. Exact union and new counts.}
For $1\leq e<T$, there are $N_T$ weights, $B_T$ positive
multiples of $268$, $A_T$ weights congruent to $-1$ modulo $67$,
and $Q_T$ weights congruent to $66$ or $200$ modulo $268$.
For a fixed $r\in\{1,\ldots,267\}$, the count of
$e\equiv r\pmod{268}$ is $\lfloor(T+267-r)/268\rfloor$.
The six residues in $\mathcal E$ are distinct and nonzero,
so exactly $R_T$ further weights are omitted from the new
$\delta_{-8}$ vanishing list.  This proves (ARW.6).

Every old pure quadratic progression $e\equiv-1\pmod{67}$
has residues $66,133,200,267$ modulo $268$.
None is zero or in $\mathcal E$.
The old trivial-label progression comprised these same residues
and the positive multiples of $268$.
Thus all old omitted triples are still present, with no overlap
between the reorganized pure and nontrivial-cubic parts.
The old four-pure-label aggregate count was $4A_T+B_T$;
the new count is $4N_T-3B_T-R_T$.
Their difference and the sum of the unchanged rows give (ARW.7).
At $e=1$, all four pure labels are newly included; consequently
new inclusion is proper for every $T\geq2$.
This is a proper reduction of the unknown term set, not an
assertion of strict numerical decrease of its nonnegative sum.

\proofstep{4. A single total-union index cost.}
Use the original dual complex
$C_{\ell,\eta,e}=R\Gamma(G_{K,\Sigma_\ell},
O_\ell(\eta^{-1})(e+1))$, with the codifferent identification
and original normalizing rank $b_\ell$ of (FL.2).
For every removed label, (EU.2) and its strict zero give
\[
 \begin{aligned}
 \frac{\log\#k_{O_\ell}}{b_\ell}\max(0,i_{\ell,\eta,e})
 &\leq\frac{\log\#H^2(C_{\ell,\eta,e})}{b_\ell}\\
 &=\frac{\left|\log\#H^2(C_{\ell,\eta,e})
                   -\log\#\mathcal S_{\ell,\eta}(e)\right|}{b_\ell}.
 \end{aligned}
 \tag{ARW.10}
\]
At each weight the complete union is one subsum of the
nonnegative all-prime error (FL.4), bounded by
$16eD_u^2\log(8eD_u)+D_ue\mathcal R_u+2D_ur_uc_0$.
Summing once over $1\leq e<T$ proves (ARW.8), using
$\sum e\leq T^2/2$ and $\log(8eD_u)\leq\log(8TD_u)$.
No extra full error is added for the broader residues.
Nonnegative additivity gives the displayed extended-sum identity
and monotonicity without subtracting infinities.

The unchanged bounds $D_u\leq288u^{2+8\varepsilon}$,
$r_u\leq12$, $\mathcal R_u=O_\varepsilon(u^{\varepsilon/5}(\log u)^2)$
and $T\leq u^{399}$ give largest powers
$802+16\varepsilon$ and $800+41\varepsilon/5$.
The full error is therefore $o(\exp(u^\varepsilon/2))$ for each
fixed positive $\varepsilon$, uniformly in the prescribed pair.
All other indices and every other coefficient prime remain.

\proofstep{5. Exactly the same common window.}
Fix the retained datum, page and $0<\varepsilon<1/256$.
Suppose the new premise supplies $J\geq1$, $X_0$ and
$s\in(0,1)$, independent of $X,q,d,u,\ell,\eta,e$, such that
for every $X\geq X_0$ and prescribed
$\operatorname{WinPair}(\varepsilon,X^{1/J},q,d)$ there is
one $W$, with $X^{1/J}\leq W$ and $W^2\leq X$, for which
\[
 I_{399}^{\rm estreflection}(H_u,u)\leq(1-s)\exp(u^\varepsilon/2)
                                  \quad(W\leq u\leq W^2).
 \tag{ARW.11}
\]
Choose $U_0\geq2$ making the uniform bound in (ARW.8) at most
$(s/2)\exp(u^\varepsilon/2)$ for every $u\geq U_0$.
Only the fixed datum, page, $\varepsilon$ and slack enter this choice.
Replace $X_0$ by $\max(X_0,U_0^J)$; the same $W$ then gives
\[
 I_{399}(H_u,u)\leq(1-s/2)\exp(u^\varepsilon/2)
                                  \quad(W\leq u\leq W^2).
 \tag{ARW.12}
\]
For a fixed finite simultaneous family use the maximum of its
uniform error constants, still on this one $W$.
No intersection of separately chosen windows is required.
Conversely either older premise implies (ARW.11) with unchanged
parameters by nonnegativity.  The same argument applies to
uniform little-$o$ premises.  Every actual field, original
coefficient, source $R.t$ and mixed-realization clause is unchanged.
\end{proof}

\paragraph{Remaining scope and geometric dictionary.}
The three nontrivial pure quadratic residue-zero rows and the
six uncovered $\delta_{-8}$ residue rows remain unknown here.
Their omission from the proved set does not assert nonvanishing.
Every other unlisted finite character and coefficient prime remains,
as do the all-prime upper bound, mixed theta/Selmer identification,
analytic realization and H04.4.
The previously audited SGA~1, Expos\'e~V, \S\S4 and 7 finite-cover
and original rank-one Artin--Tate dictionary is retained.
No new motive, cohomology comparison, globally maximal vanishing
classification or Lean execution is asserted.
